% Sub-stage F2: the stages of the sieve and the check: the reductions by a multiplication with the
% reciprocal of a modulus, the test for common factors by the binary algorithm, the third stage
% numerator by numerator, the second stage on the bit arrays that the first leaves nonzero, the
% staged sieve with the cut points of the parameters, and the parameters of the command line by the
% counting rules of ratpoints. Proposed names are set with \lname and become \lean once the
% declarations exist.
\chapter{The stages of the sieve and the check}\label{chap:stages}

F2 models the three stages in which ratpoints applies its moduli to the numerators of a
denominator, and the check of a candidate. The first stage and-s the rows of the first $s_1$ moduli
into every bit array of a piece, as in Chapter~\ref{chap:words}; the second and-s the rows of the
moduli from $s_1$ to $s_2$ into each bit array that the first stage leaves nonzero, one after the
other, and stops when no bit is left; and each numerator whose bit survives goes to the test for
common factors with the denominator, by the binary algorithm of \verb|relprime|, and then to the
third stage, which looks up, for each prime from $s_2$ on, whether the form can be a square at that
numerator, before the exact check. The reductions of the third stage and of the Jacobi test multiply
by the reciprocal of the modulus, as the C code does; the cut points $s_1$ and $s_2$ come from the
parameters, and the parameters of the command line take them by the counting rules of ratpoints.

Each stage only tests the moduli of the parameters, so the numerators that reach the exact check are
those of the unstaged sieve of Chapter~\ref{chap:words} with its check, in the same order, and the
theorems of Chapter~\ref{chap:denominators} keep their statements. The statements that comparator
certifies keep their text too: the well-formedness of the parameters (Definition~\ref{def:Params})
has required the moduli of the third stage to be primes of the table since stage~F1, which is what
the third stage needs. What changes for comparator is what \verb|wf_sieveParams| states, through
\verb|sieveParams|, which the challenge imports from the model: the parameters of the command line
now take their three stages by the counting rules of ratpoints, and comparator checks the statement
again; and the proofs of the theorems change through the contracts of the run, of
\verb|sieving_info| and of the Jacobi test on which they rest.

\section{What the C code does in the stages}\label{sec:stagesC}

For a denominator $b$, \verb|sift| (\csource{find_points.c}{3442}{3617}) sets the rows of $b$ for
the moduli of the first two stages, clears the flag \verb|stage3_filled| of the sieving thread,
whether the table of the third stage has been filled for $b$, and sets its flag \verb|compute_bc|,
that the products of the check are still to be computed. Then for each piece of its numerators,
\verb|_ratpoints_sift0| (\csource{sift.c}{498}{1305}):
\begin{enumerate}
  \item and-s the rows of the moduli $0, \dots, s_1 - 1$ into every bit array of the piece, from the
    pattern of the 2-adic class, and applies the masks of the ends of the piece;
  \item scans the bit arrays from the first; for each nonzero one with the word number $i$, copies
    it and and-s into the copy the bit array $(i + \mathrm{offset}) \bmod p$ of the row of each
    modulus $p$ from $s_1$ to $s_2 - 1$ in turn, stopping when the copy is zero
    (\csource{sift.c}{1174}{1202}); with \verb|USE_LONG_IN_PHASE_2|, a build option that the default
    build leaves off, it does this one word of the copy at a time (\csource{sift.c}{1113}{1173});
  \item for each set bit of what is left, lowest first, tests the numerator $a$ it stands for:
    \verb|accepted| (\csource{sift.c}{438}{446}) runs \verb|relprime| (\csource{sift.c}{318}{343}),
    the binary algorithm for $\gcd(a, b) = 1$ with a loop body without branches, then fills the
    table of the third stage if it is not filled yet (\verb|fill_checks|,
    \csource{sift.c}{405}{431}) and runs \verb|stage3| (\csource{sift.c}{291}{312}), which for each
    prime $p$ of the third stage that does not divide $b$ rejects $a$ when $f(a b^{-1})$ is not a
    square modulo $p$; a numerator that passes goes to the exact check, which computes the products
    $c_k b^{d - k}$ when the first numerator of $b$ reaches it
    (\csource{find_points.c}{5250}{5338}).
\end{enumerate}
The table of the third stage holds, for each prime $p$: the inverse of $b$ modulo $p$, $0$ when
$p \mid b$, from the residue $b \bmod p$ computed by \verb|RP_MULMOD| for $b < 2^{32}$ and by
\verb|%| above; the reciprocal of $p$; and the bias $p H$ when $2 p H < 2^{32}$, else $0$. With the
bias, \verb|stage3| reduces the value $a b^{-1} + p H$, which lies in $[H, 2 p H - H]$ for
$|a| \le H$, by \verb|RP_MULMOD|; without it, it divides. \verb|RP_MULMOD| and \verb|RP_MULDIV|
(\csource{rp-private.h}{404}{417}) compute the remainder and the quotient of $u$ by $p$ from the
reciprocal $m = \lfloor (2^{64} - 1) / p \rfloor + 1$ with a 128-bit product; the C code relies on
$u < 2^{32}$, and they are exact whenever $u p < 2^{64}$. The Jacobi test by tables
(Definition~\ref{def:jacobiTables}) divides by its primes with \verb|RP_MULDIV|.

\verb|sieving_info| chooses the cut points: $s_1$ moduli for the first stage, up to $s_2$ for the
second, and the primes of the third stage up to $s_3$, by the options \verb|-n|, \verb|-N| and
\verb|-P| where they are given and by the cost of each stage, which depends on the curve, where they
are not (\csource{find_points.c}{2803}{3034}). Its third stage takes primes only, and only those
that exclude a residue.

Why the stages change nothing: the and of the rows is commutative and associative, and a bit array
that is zero stays zero, so the second stage leaves exactly the bits that the first stage would
leave with the moduli of both; for a numerator $a$ prime to $b$, the third stage passes $a$ exactly
when the row of $b$ of every prime of the stage admits $a$, since for $p \nmid b$ the row of $b$
reads the same table at the same residue $a b^{-1} \bmod p$ (Definition~\ref{def:admits}), and for
$p \mid b$ the row admits every numerator prime to $p$; and a numerator with a common factor has no
point. So the numerators that reach the exact check are those that survive all the moduli and are
prime to $b$: those that the unstaged sieve keeps and that the test of the $\gcd$ in its check
(Definition~\ref{def:checkNum}) lets through. The stages, like the caches, are a data refinement of
the unstaged sieve with its check: the same numerators reach the exact check, in the same order.

\section{The reductions by a multiplication}\label{sec:mulMod}

\begin{definition}[The reciprocal and the reductions]\label{def:mulMod}
  Proposed names: \lname{Ratpoints.stage3Limit}, \lname{Ratpoints.stage3Limit\_le\_two\_pow\_32},
  \lname{Ratpoints.Sieve.magic}, \lname{Ratpoints.Sieve.mulMod}, \lname{Ratpoints.Sieve.mulDiv}.
  The bound $\mathrm{stage3Limit} = 2^{32}$ below which the third stage reduces by a multiplication
  (\verb|RP_STAGE3_LIMIT|, \csource{rp-private.h}{472}{472}), with the fact
  $\mathrm{stage3Limit} \le 2^{32}$ checked by evaluation; the reciprocal of a modulus $p$,
  $\mathrm{magic}(p) = \lfloor (2^{64} - 1) / p \rfloor + 1$, computed where it is read, as
  \verb|sieving_info| collects those of the moduli (\csource{find_points.c}{3064}{3070}); the
  remainder of $u$ by $p$ from the reciprocal $m$, the low $64$ bits of $m u$, times $p$, shifted
  right by $64$ bits (\csource{rp-private.h}{404}{408}, \csource{rp-private.h}{411}{411},
  \csource{rp-private.h}{413}{414}, \csource{rp-private.h}{417}{417}); and the quotient, the high
  $64$ bits of $u m$ (\csource{rp-private.h}{409}{410}, \csource{rp-private.h}{415}{416}): the
  forms with a 128-bit product, which the default build and the build of the model compile, where
  the others divide.
\end{definition}

\begin{lemma}[The reductions are exact]\label{lem:mulMod}
  \uses{def:mulMod}
  Proposed names: \lname{Ratpoints.Sieve.magic\_mul\_eq}, \lname{Ratpoints.Sieve.mulMod\_magic},
  \lname{Ratpoints.Sieve.mulDiv\_magic}, \lname{Ratpoints.Sieve.sub\_mulDiv\_magic\_mul},
  \lname{Ratpoints.mul\_lt\_two\_pow\_64\_of\_le\_mulModLimit}. For $0 < p$ and $u p < 2^{64}$, the
  reduction with the reciprocal of $p$ gives $u \bmod p$, and the quotient gives $\lfloor u / p
  \rfloor$.
\end{lemma}

\begin{proof}
  Write $\mathrm{magic}(p)\, p = 2^{64} + \delta$ with $0 \le \delta < p$ and $u = k p + s$ with
  $s < p$. Then $\mathrm{magic}(p)\, u = k\, 2^{64} + (s\, 2^{64} + u \delta) / p$, whose second
  part is an integer below $2^{64}$, since $s < p$ and $u \delta \le u p < 2^{64}$; so the low half
  is $(s\, 2^{64} + u \delta)/p$ and the high half $k$. The low half times $p$ is
  $s\, 2^{64} + u \delta$, whose high half is $s$.
\end{proof}

The callers satisfy the bound: the denominators that \verb|fill_checks| reduces with the reciprocal
(it takes the remainder of the larger ones) and those of the Jacobi test are at most
$\mathrm{mulModLimit} = 2^{32} - 1$, below $2^{32}$ by a fact checked by evaluation
(\lname{Ratpoints.mulModLimit\_lt\_two\_pow\_32}, in the node of the Jacobi test by tables), and
their moduli are below $2^{10}$; the values that \verb|stage3| reduces are below $2 p H <
\mathrm{stage3Limit} \le 2^{32}$. The Jacobi test by tables (Definition~\ref{def:jacobiTables})
takes the quotient and the remainder by each of its primes by the quotient with the reciprocal, as
\verb|jacobi_test| does; its contract takes $b q < 2^{64}$ for the primes $q$ of the Jacobi
information, which the set-up of the test now gives as primes of the table, and the statement of the
test for the denominators up to $b_{\mathrm{high}} \le \mathrm{mulModLimit}$ stays as it is
(\lname{Ratpoints.Sieve.jacobiSetup\_spec}, \lname{Ratpoints.Sieve.jacobiTest\_spec},
\lname{Ratpoints.Sieve.run\_jacobiTest}).

\section{The test for common factors}\label{sec:relprime}

\begin{definition}[The binary algorithm of \texttt{relprime}]\label{def:relprime}
  Proposed names: \lname{Ratpoints.Sieve.ctzNat}, \lname{Ratpoints.Sieve.relprimeOdd},
  \lname{Ratpoints.Sieve.relprime}.
  The number of trailing zeros of a natural number $n \ge 1$, by the loop of \verb|RP_CTZL|, which
  halves $n$ while it is even (\csource{rp-private.h}{429}{444}); the loop is modelled twice, on
  words in Definition~\ref{def:ctz} and on natural numbers here. For odd $m, n \ge 1$: while
  $n \ne m$, $n$ becomes $\min(m, n)$ and $m$ the odd part of $|m - n|$; the result is whether $m$
  ends as $1$ (\csource{sift.c}{334}{342}). For an integer $a$ and $b \ge 1$: for $a = 0$, whether
  $b = 1$; otherwise, for $m = |a|$, false when $m$ and $b$ are both even, and else the loop on the
  odd parts of $m$ and $b$ (\csource{sift.c}{318}{332}, \csource{sift.c}{343}{343}).
\end{definition}

\begin{lemma}[The test for common factors]\label{lem:relprime}
  \uses{def:relprime}
  Proposed names: \lname{Ratpoints.Sieve.CtzNatInvariant}, \lname{Ratpoints.Sieve.ctzNat\_spec},
  \lname{Ratpoints.Sieve.RelprimeOddInvariant}, \lname{Ratpoints.Sieve.relprimeOdd\_spec},
  \lname{Ratpoints.Sieve.relprime\_spec}, \lname{Ratpoints.fitsLong\_natCast\_sub\_natCast}. For $n
  \ge 1$ the number of trailing zeros $t$ has $n = 2^t n'$ with $n'$ odd; for odd $m, n \ge 1$ the
  loop gives whether $\gcd(m, n) = 1$; and for $b \ge 1$ the test gives whether $\gcd(a, b) = 1$.
  The values of the loop stay at most $\max(m, n)$, so that the \verb|long|s of the C code stay in
  range for $|a|, b < 2^{63}$.
\end{lemma}

\begin{proof}
  Halving keeps $n = 2^t n'$ and ends since $n$ decreases. In the loop, both numbers stay odd and
  their $\gcd$ stays the same: $\gcd(m, n) = \gcd(\min(m, n), |m - n|)$, and the power of $2$ in
  $|m - n|$ is prime to the odd $\min(m, n)$; the sum $m + n$ decreases, since
  $\min(m, n) + |m - n| / 2 < m + n$; at the end $m = n = \gcd$. For $a \ne 0$: when $m$ and $b$ are
  both even, $2$ divides both; otherwise removing the powers of $2$ does not change whether the
  $\gcd$ is $1$, since one of the two is odd.
\end{proof}

\section{The third stage}\label{sec:stage3}

\begin{definition}[The third stage]\label{def:stage3}
  \uses{def:mulMod, def:SieveEntry, def:jacobiTables}
  Proposed names: \lname{Ratpoints.Sieve.CheckSpec}, \lname{Ratpoints.Sieve.checkOf},
  \lname{Ratpoints.Sieve.fillChecks}, \lname{Ratpoints.Sieve.CheckSpec.index},
  \lname{Ratpoints.Sieve.stage3}, \lname{Ratpoints.Sieve.CheckSpec.admits},
  \lname{Ratpoints.Sieve.SieveEntry.IsPrimeEntry}. What the third stage needs of a prime $p$ for the
  denominator $b$ (\csource{rp-private.h}{465}{466}): $p$, the inverse $b^{-1}$ modulo $p$ read from
  the table of inverses of $p$ at $b \bmod p$, computed by the reduction with the reciprocal for $b
  \le \mathrm{mulModLimit}$ and by the remainder above, or $0$ when $p \mid b$; the bias $p H$ when
  $2 p H < \mathrm{stage3Limit}$ and $0$ otherwise; the reciprocal of $p$; and the table of $p$
  (\csource{sift.c}{414}{428}). For a sieve entry of another shape than a prime with its table,
  which the third stage of well-formed parameters never holds, the inverse is $0$, so that the stage
  skips it. The list of these for the sieve entries of the third stage, in their order
  (\csource{sift.c}{405}{413}, \csource{sift.c}{429}{431}). The residue at which the table of $p$ is
  read for the numerator $a$: with the bias, $a b^{-1} + p H$ reduced with the reciprocal; without
  it, the remainder of $a$ by the truncating division, plus $p$ if it is negative, times $b^{-1}$,
  reduced by the remainder (\csource{sift.c}{302}{308}). The third stage for $a$: for each prime in
  turn, skipped when the inverse is $0$, $a$ is rejected when the table says that $f$ is not a
  square modulo $p$ at that residue (\csource{sift.c}{291}{301}, \csource{sift.c}{309}{312}). A
  \emph{prime entry} is a sieve entry whose single factor is a prime $1 < p \le 2^{32}$ with its
  table, which holds for each nonzero residue an inverse in $[1, p - 1]$.
\end{definition}

\begin{lemma}[The third stage tests the rows]\label{lem:stage3}
  \uses{def:stage3, lem:mulMod, def:admits, def:SieveEntry, def:examinePrime}
  Proposed names: \lname{Ratpoints.Sieve.Stage3Invariant}, \lname{Ratpoints.Sieve.stage3\_spec},
  \lname{Ratpoints.Sieve.CheckSpec.index\_eq}, \lname{Ratpoints.Sieve.le\_mul\_add\_mul\_and\_le},
  \lname{Ratpoints.Sieve.toNat\_emod\_mul\_lt\_sq}, \lname{Ratpoints.Sieve.binv\_checkOf},
  \lname{Ratpoints.Sieve.CheckSpec.admits\_checkOf},
  \lname{Ratpoints.Sieve.run\_stage3\_fillChecks},
  \lname{Ratpoints.Sieve.getElem!\_inversesMod\_lt},
  \lname{Ratpoints.Sieve.isPrimeEntry\_examinePrime}. For prime entries, $|a| \le H$ and $\gcd(a, b)
  = 1$: the residue at which the table of $p$ is read is $a b^{-1} \bmod p$ for $p \nmid b$, in both
  branches of the bias and on both sides of $\mathrm{mulModLimit}$, and the third stage passes $a$
  exactly when every entry's row of $b$ admits $a$. The values that the C code holds in \verb|long|s
  stay in range: $a b^{-1} + p H$ lies in $[H, 2 p H - H]$, and the product of the residue and the
  inverse is below $p^2$. The entry of a prime $p \le 2^{32}$ with its table
  (Definition~\ref{def:examinePrime}) is a prime entry.
\end{lemma}

\begin{proof}
  \uses{lem:mulMod}
  A prime that does not divide $b$ is not skipped: the inverse that its table holds for a nonzero
  residue lies in $[1, p - 1]$, which also puts $a b^{-1} + p H$ in $[H, 2 p H - H]$ and the product
  of the residue and the inverse below $p^2$. The reductions are exact by Lemma~\ref{lem:mulMod}: $b
  p < 2^{64}$ for $b \le \mathrm{mulModLimit} < 2^{32}$ and $p \le 2^{32}$, and the biased value is
  below $2 p H < 2^{32}$. The biased value is congruent to $a b^{-1}$, and the truncated remainder
  plus $p$ is congruent to $a$. A prime that divides $b$ is skipped, and its row admits $a$ since $p
  \nmid a$ by the $\gcd$; for the others, the row of $b$ admits $a$ exactly when the table of $p$ is
  true at $a b^{-1} \bmod p$ (Definition~\ref{def:admits}). The inverses of the table of a prime are
  those of the nonzero residues.
\end{proof}

\begin{definition}[The test of a candidate and the exact check]\label{def:accepted}
  \uses{def:relprime, def:stage3, def:bruteForce, def:swapPoint}
  Proposed names: \lname{Ratpoints.Sieve.accepted}, \lname{Ratpoints.Sieve.checkPoint}.
  A numerator $a$ goes on to the exact check for the denominator $b$ when it is prime to $b$
  (Definition~\ref{def:relprime}) and passes the third stage, whose list is a thunk made once per
  denominator and forced when the first numerator of $b$ gets that far, for \verb|stage3_filled|
  (\csource{sift.c}{438}{446}). The exact check of $a$ from the products $c_k b^{d - k}$, a thunk
  made once per denominator for \verb|compute_bc|, and the factor $b$ for odd $d$: the points
  $(a, y, b)$ and, for $y \ne 0$, $(a, -y, b)$, each mapped by $\mathrm{swapPoint}$ when the flag of
  the reversal is set (\csource{sift.c}{202}{205}, \csource{find_points.c}{3122}{3122},
  \csource{find_points.c}{5250}{5338}, \csource{rp-private.h}{596}{597}); the check of a numerator
  of Definition~\ref{def:checkNum}, which tests the $\gcd$ itself, becomes the reference of the two.
\end{definition}

\begin{lemma}[The test of a candidate]\label{lem:accepted}
  \uses{def:accepted, lem:relprime, lem:stage3, def:checkNum}
  Proposed names: \lname{Ratpoints.Sieve.accepted\_spec}, \lname{Ratpoints.Sieve.run\_accepted},
  \lname{Ratpoints.Sieve.checkPoint\_eq\_checkNum},
  \lname{Ratpoints.Sieve.ite\_run\_accepted\_checkPoint}. For $b \ge 1$, $|a| \le H$ and prime
  entries in the third stage, with no hypothesis on $\gcd(a, b)$, a numerator passes the test
  exactly when it is prime to $b$ and every entry's row of $b$ admits it; the exact check of a
  numerator prime to $b$ gives the points of the reference check.
\end{lemma}

\begin{proof}
  Lemma~\ref{lem:relprime} for the first test, which rejects every numerator with a common factor;
  Lemma~\ref{lem:stage3} for the second, which runs only after the first has passed, so that its
  hypothesis $\gcd(a, b) = 1$ holds; a thunk has the value of its function.
\end{proof}

\section{The second stage}\label{sec:secondStage}

\begin{definition}[The second stage]\label{def:secondStage}
  \uses{def:bitArrayOps, def:rows}
  Proposed names: \lname{Ratpoints.useLongInPhase2}, \lname{Ratpoints.Sieve.secondStage},
  \lname{Ratpoints.Sieve.secondStageWord}. Whether the second stage runs one word at a time,
  $\mathrm{useLongInPhase2} = \mathrm{false}$ as in the default build (\verb|USE_LONG_IN_PHASE_2|,
  \csource{sift.c}{1084}{1084}, \csource{sift.c}{1086}{1086}, \csource{sift.c}{1113}{1113},
  \csource{sift.c}{1174}{1174}, \csource{sift.c}{1297}{1297}). The second stage on a copy of a bit
  array with the word number $i$: for each modulus $p$ with its row in turn, while the copy is
  nonzero, the copy is and-ed with bit array $i \bmod p$ of the row (\csource{sift.c}{1083}{1083},
  \csource{sift.c}{1085}{1085}, \csource{sift.c}{1175}{1202}); and on one word $\kappa$ of it,
  likewise with word $\kappa$ of those bit arrays (\csource{sift.c}{1135}{1136},
  \csource{sift.c}{1141}{1152}).
\end{definition}

\begin{lemma}[The second stage]\label{lem:secondStage}
  \uses{def:secondStage}
  Proposed names: \lname{Ratpoints.Sieve.SecondStageInvariant},
  \lname{Ratpoints.Sieve.secondStage\_spec}, \lname{Ratpoints.Sieve.secondStageWord\_spec},
  \lname{Ratpoints.Sieve.andWords\_append}, \lname{Ratpoints.Sieve.andWords\_rbaPack\_mul\_add}.
  The second stage gives the copy and-ed with the bit arrays $i \bmod p$ of the rows of all its
  moduli, and on one word the same word of that; and-ing the rows of a list of moduli and then those
  of a second list is and-ing those of the concatenation.
\end{lemma}

\begin{proof}
  An invariant of the loop over the moduli: the copy is the and of the rows so far, or zero, and
  then the and of all of them is zero too.
\end{proof}

\begin{definition}[A bit array through the second stage and the extraction]\label{def:siftBitArray}
  \uses{def:secondStage, def:walkWord, def:accepted}
  Proposed names: \lname{Ratpoints.Sieve.siftBitArray}.
  For a bit array that the first stage leaves nonzero, with the word number $i$, from a copy of it:
  in the arm that an argument chooses, $\mathrm{useLongInPhase2}$ in the run, either the second
  stage on the whole copy and, if it is still nonzero, the set bits of its words walked through in
  turn, or for each nonzero word the second stage on that word and the walk through its set bits,
  bit $t$ of word $\kappa$ standing for the numerator $R i + 64 \kappa + t$; each numerator goes to
  the test of a candidate and, if it passes, to the exact check (\csource{sift.c}{1094}{1111},
  \csource{sift.c}{1126}{1134}, \csource{sift.c}{1137}{1139}, \csource{sift.c}{1154}{1159},
  \csource{sift.c}{1172}{1173}, \csource{sift.c}{1204}{1210}, \csource{sift.c}{1213}{1226},
  \csource{sift.c}{1257}{1264}, \csource{sift.c}{1293}{1296}; with \verb|EXT0| and \verb|EXT|,
  \csource{rp-private.h}{135}{136}, \csource{rp-private.h}{177}{178},
  \csource{rp-private.h}{238}{239}, \csource{rp-private.h}{264}{266},
  \csource{rp-private.h}{288}{289}).
\end{definition}

\begin{lemma}[A bit array through the second stage and the extraction]\label{lem:siftBitArray}
  \uses{def:siftBitArray, lem:secondStage, def:walkWord}
  Proposed names: \lname{Ratpoints.Sieve.SiftBitArrayInvariant},
  \lname{Ratpoints.Sieve.SiftBitArrayWalkInvariant}, \lname{Ratpoints.Sieve.siftBitArray\_spec}.
  In both arms, the points appended are those of the numerators of the bit array whose bits
  survive the and of the copy with the rows of the moduli of the second stage and that pass the
  test of a candidate, in increasing order.
\end{lemma}

\begin{proof}
  Lemma~\ref{lem:secondStage} and the contract of the walk (Definition~\ref{def:walkWord}), in each
  arm with its own invariant: a word that is zero, or a bit array that the second stage leaves zero,
  has no set bits to walk.
\end{proof}

\section{The staged sieve}\label{sec:staged}

The walk through the set bits of a word (Definition~\ref{def:walkWord}) runs the test of a candidate
and then the exact check for each set bit; the scan (Definition~\ref{def:scanPiece}) copies each
nonzero bit array and hands it to Definition~\ref{def:siftBitArray}, with the moduli of the second
stage; one piece (Definition~\ref{def:siftPiece}) runs the first stage with the first $s_1$ moduli
of its list and the scan with the rest; and one denominator (Definition~\ref{def:sift}) takes the
end $s_1$ of the first stage, the rows of the moduli of the first two stages and the sieve entries
of the third, and makes the thunks of the products and of the table of the third stage. Their
contracts keep their form: the walk appends the points of the numerators of its set bits that pass
the test of a candidate, the scan those of the numerators whose bits survive the first stage and the
moduli of the second and that pass the test of a candidate, and the piece those of the numerators
that survive the first stage, the masks and the second stage and pass the test. The lemma that the
words compute the Booleans (Lemma~\ref{lem:sift_eq_checkDen}) keeps its conclusion, the points of
the sieve of stage~C1 over all the sieve entries followed by its check, and gains the hypotheses
that $b \ge 1$ and that the entries of the third stage are prime entries. It rests on the following.

\begin{lemma}[The stages compute the unstaged sieve]\label{lem:staged}
  \uses{lem:secondStage, lem:accepted, lem:siftBitArray, def:sift}
  Proposed names: \lname{Ratpoints.Sieve.acceptedPoints},
  \lname{Ratpoints.Sieve.acceptedPoints\_eq\_ite}, \lname{Ratpoints.Sieve.survives\_append}. The
  points that the extraction gives for a numerator $a$ with $|a| \le H$, those of the exact check
  when $a$ passes the test of a candidate and none otherwise, are those of the reference check when
  $a$ survives the entries of the third stage, and none otherwise; a numerator survives a list of
  entries followed by another when it survives both. So for $s_1 \le s_2 \le$ the number of moduli,
  the numerators whose bits survive the first $s_1$ moduli in the first stage and the next $s_2 -
  s_1$ in the second, and that pass the test of a candidate with the entries from $s_2$ on, give the
  points that the reference check gives for the numerators that survive all the entries.
\end{lemma}

\begin{proof}
  Lemma~\ref{lem:accepted} for the first part, since the reference check gives no points for a
  numerator with a common factor; Lemma~\ref{lem:secondStage} for the split of the list at $s_1$.
\end{proof}

The run over the denominators (Definition~\ref{def:sieveRun}) takes the cut points $s_1$ and $s_2$:
it makes the rows of the first $s_2$ sieve entries, and hands the rest to \verb|sift| for the third
stage, for the ends of the second and the third stage, which the run takes from the arguments as its
first decision (\csource{find_points.c}{3861}{3861}) and \verb|run_block| sets for each block
(\csource{find_points.c}{3911}{3911}). Its contract keeps its conclusion and needs that the
denominators are positive, $b_{\mathrm{low}} \ge 1$, since the test for common factors does not end
for $b = 0$ and an odd numerator, and that the entries from $s_2$ on are prime entries; the search
(Definition~\ref{def:searchPoints}) passes $b_{\mathrm{low}} = 1$ and the cut points of its
parameters, and what \verb|sieving_info| computes (Definition~\ref{def:sievingInfo}) gives entries
of that shape for the first $\mathrm{numPrimes}$ primes of the table, one for each modulus in its
order, which the well-formedness makes the moduli of the third stage
(\lname{Ratpoints.makeModulus?\_spec},
\lname{Ratpoints.factors\_eq\_of\_run\_makeModulus?\_eq\_some}, \lname{Ratpoints.sievingInfo\_spec},
\lname{Ratpoints.Params.WF.exists\_getElem!\_primeTable\_eq\_of\_mem\_drop}).

\section{The cut points of the command line}\label{sec:cuts}

The parameters of the executable (Definition~\ref{def:sieveParams}) take the cut points by the
counting rules of ratpoints, where they do not depend on the curve. For the number
$q = \min(k, \mathrm{numPrimes})$ of primes examined, with the value $k \ge 0$ of \verb|-p| or else
$\mathrm{defaultNumPrimes}$, the bound $L$ on the moduli is $q$ under \verb|-p|, which is a hard
limit for every stage (\csource{find_points.c}{4577}{4577}), and $\mathrm{numPrimes}$ without it,
since ratpoints then looks at further primes when its stages want more
(\csource{find_points.c}{2693}{2705}, \csource{find_points.c}{2980}{2980}). With the values $n$, $N$
and $P$ of \verb|-n|, \verb|-N| and \verb|-P|, each absent when negative, and the constants
\lname{Ratpoints.defaultSp1} $= 11$ (\verb|RATPOINTS_DEFAULT_SP1|) and \lname{Ratpoints.sp2Extra}
$= 11$ (\verb|RATPOINTS_SP2_EXTRA|), which join the node of the parameters of the executable:
\begin{itemize}
  \item $N$ is capped at $q$, and the first stage takes $s_1 = \min(n, N)$ moduli, or
    $\min(n, L)$ without \verb|-N|, or $\min(\mathrm{defaultSp1}, L)$ without \verb|-n|, a count
    that an option gives being taken as it is (\csource{ratpoints.h}{42}{42},
    \csource{find_points.c}{2503}{2503}, \csource{find_points.c}{2845}{2846},
    \csource{find_points.c}{4655}{4656});
  \item the second stage ends at $s_2 = \max(s_1, N)$, the first two stages topped up to $N$, or
    without \verb|-N| at $\min(s_1 + \mathrm{sp2Extra}, L)$, the offset that ratpoints adds unscaled
    under \verb|-U 0| (\csource{ratpoints.h}{95}{97}, \csource{find_points.c}{1619}{1623},
    \csource{find_points.c}{2571}{2572}, \csource{find_points.c}{2599}{2599},
    \csource{find_points.c}{2808}{2808}, \csource{find_points.c}{2815}{2815},
    \csource{find_points.c}{2864}{2869});
  \item the moduli are the first $\min(s_2 + P, L)$ primes of the table, or without \verb|-P| the
    first $\max(s_2, q)$, so that the third stage holds the rest of the first $q$
    (\csource{find_points.c}{2922}{2922}, \csource{find_points.c}{2926}{2928},
    \csource{find_points.c}{2949}{2949}, \csource{find_points.c}{2973}{2975},
    \csource{find_points.c}{3015}{3015}, \csource{find_points.c}{3017}{3017}).
\end{itemize}
A run without options has $11$ moduli in the first stage, $11$ more in the second and $8$ in the
third, the first $30$ primes in all, and uses all three stages. The parameters are well-formed as
before (\verb|wf_sieveParams|, whose text is unchanged): the moduli are a prefix of the first
$\mathrm{numPrimes}$ primes of the table, and $s_1 \le s_2 \le$ their number. Where ratpoints
decides by the curve, it may take fewer moduli in a stage: the primes that say nothing leave its
pool, those that say only that the numerator is prime to the denominator stay in the pool of the
first two stages without counting toward the primes it looks for and leave that of the third, and
its composite moduli take several primes of the pool. With \verb|-n|, \verb|-N| (or \verb|-U 0|) and
\verb|-P| given, the first stage, the first two stages and the moduli in all are at most these
counts; where an option is absent, ratpoints chooses $s_1$, the offset of the second stage and the
length of the third by their costs, its third stage taking the primes of its pool in increasing
order of their densities, which may give more moduli in a stage or fewer. The cut points change
neither the points nor their order, so the model takes the counts that the options give and states
the rest as a deviation.

\section{The command-line program and the differential tests}\label{sec:stagesCli}

The program runs the three stages with the cut points of its parameters; without options, the moduli
of the first $30$ primes are now split $11$, $11$ and $8$, where they were all in the first stage.
The differential tests compare every point list in its order, as before; the suite of the sieving
parameters (Section~\ref{sec:sieveCli}) gains option sets of \verb|-n|, \verb|-N| and \verb|-P| that
run the second and the third stage, at heights at which a run with few moduli before the third stage
stays short, since then most numerators reach the test of a candidate, which the model runs
numerator by numerator, as the C code does but more slowly. \verb|#guard| examples compare each new
component with values of the C program and with the reference: the reductions with the remainder and
the quotient at the ends of their range, the test for common factors with the $\gcd$ beyond
$2^{64}$, the third stage after its table with the rows of its primes in both branches of the bias
and above $\mathrm{mulModLimit}$, both arms of the second stage with the and of the rows, the piece
and the denominator at many cut points with the unstaged sieve, and the stage sizes of the
parameters with those that \verb|ratpoints -v| reports, with \verb|-U 0| where \verb|-N| is absent.

\section{Deviations}\label{sec:stagesDeviations}

Each component of this chapter differs from the C code in a stated way, recorded in its
\verb|Deviation:| paragraph:
\begin{itemize}
  \item natural numbers and integers for the \verb|long|s and \verb|unsigned long|s: the reductions
    compute the exact values of the 128-bit products, the test for common factors runs on natural
    numbers, so that it ends for every height, with the minimum and the absolute difference for the
    loop body without branches and the mask of an arithmetic shift, and its trailing zeros by the
    loop of \verb|RP_CTZL| on natural numbers; the bounds that keep the C values in range are lemmas
    (Lemmas~\ref{lem:relprime} and~\ref{lem:stage3});
  \item the reciprocal computed where it is read, for the field of the sieve entry, the array
    \verb|magics| that \verb|sieving_info| fills and the field \verb|magic| of the tables of the
    Jacobi test; the comparison $2 p H < \mathrm{stage3Limit}$ of natural numbers, for the
    comparison of \verb|double|s, which has the same outcome for every $H < 2^{63}$;
  \item the table of the third stage and the products of the check as thunks made once per
    denominator, for the arrays of the sieving thread and the flags \verb|stage3_filled| and
    \verb|compute_bc|; the table a list, with no bound $\mathrm{RATPOINTS\_NUM\_PRIMES}$ on its
    length, which ratpoints never reaches since each of its sieve entries owns a prime, where
    well-formed parameters may repeat a prime; an entry of another shape skipped;
  \item the arm of the second stage an argument, so that the guards run both in one build; the row
    of the word number $i$ at index $i \bmod p$, where ratpoints adds the offset of the packing and
    reduces by a multiplication where that is exact (stage~F3);
  \item the cut points of the parameters for the whole run, where \verb|run_block| sets the end of
    the second and the third stage for each block from the correction of the moduli (stage~I); the
    parameters of the command line by the counts that the options give, with the flat offset of the
    second stage, for the choices by cost, the scaled offset and \verb|-R|, which the model ignores;
  \item no counts of the numerators that reach the test for common factors and the third stage,
    which ratpoints reads to correct the moduli between the blocks (stage~I), and no counts, stops
    and output of the instrumentation of the stages.
\end{itemize}
The lines of this part of the C code that are neither modelled nor left to a later stage are
excluded in \verb|correspondence/ratpoints.toml|: the instrumentation of \verb|sift.c| between
statements that are not modelled, with the arms of its production build; the inlining attributes
of \verb|rp-private.h|, which carry no meaning; and \verb|RATPOINTS_DEFAULT_SP2|, which the C code
reads nowhere. Left to stage~F3 are the reductions of the row indices of the first two stages
(\verb|mod_mul|, the second copy of \verb|mod| in \verb|sift.c|, the test \verb|small|) with the
offsets of the packing and the residues of $b\, 2^{-k}$ (\verb|fill_bp_list|); left to stage~I the
correction of the moduli between the blocks, the choice of the cut points by cost, and, with the
blocks, the reduction of the first word of a block modulo the forbidden divisors tested with words,
which still divides (\verb|wordBits|).

\section{What changes in the earlier chapters}\label{sec:stagesEarlier}

Each item of the order of work changes the nodes of Chapters~\ref{chap:sieve}
to~\ref{chap:tables} that its declarations belong to, and the nodes whose C lines its declarations
take, in the same commit, so that the names of the nodes exist and their C links are the ranges of
their declarations:
\begin{itemize}
  \item the introduction of Chapter~\ref{chap:sieve} counts the stages of the sieve and the check of
    stage~F2 among the data refinements, with the bit arrays of stage~C2, the faithful tables of
    stage~F1 and the packing of stage~F3, and no longer counts the third stage among the sound
    filters: it tests the rows of moduli that the unstaged sieve applies already; the list of
    deviations of Chapter~\ref{chap:sieve} keeps, of its item on the moduli applied before the test
    for common factors, only the packing (stage~F3) and the correction between the blocks
    (stage~I), and loses its items on the plain $\gcd$ and on the products computed for every
    denominator;
  \item the parameters of the executable (Definition~\ref{def:sieveParams}) take the cut points of
    Section~\ref{sec:cuts}, with the constants $\mathrm{defaultSp1}$ and $\mathrm{sp2Extra}$; the
    paragraph after them says that the cut points are read since stage~F2 and that the stages of
    ratpoints may be longer or shorter where it decides by the curve;
  \item the first stage (Definition~\ref{def:andRows}) applies the first $s_1$ moduli only, and
    gives the lines of the second and the third stage to Definitions~\ref{def:secondStage}
    and~\ref{def:stage3}; the check of a numerator (Definition~\ref{def:checkNum}) becomes the
    reference, with no C links, its lines going to Definitions~\ref{def:relprime}
    and~\ref{def:accepted}; the walk, the scan, the piece and one denominator
    (Definitions~\ref{def:walkWord} to~\ref{def:sift}) take the test of a candidate, the second
    stage, the cut points and the thunks: the walk gains the walk of the arm one word at a time
    (\csource{sift.c}{1161}{1171}) and uses Definitions~\ref{def:ctz} and~\ref{def:accepted} instead
    of Definition~\ref{def:checkNum}; the scan copies the bit array, keeps
    \csource{sift.c}{1061}{1082}, \csource{sift.c}{1088}{1092}, \csource{sift.c}{1298}{1299} and the
    lines of \verb|zero|, gives the extraction and the lines of \verb|EXT0| and \verb|EXT| to
    Definition~\ref{def:siftBitArray} and line~1297 to Definition~\ref{def:secondStage}, and uses
    Definition~\ref{def:siftBitArray} instead of Definition~\ref{def:walkWord}; the piece gains the
    prototype of \verb|_ratpoints_sift0| (\csource{rp-private.h}{610}{614}); one denominator gains
    the set-up of the sieving thread for the cut points, the table of the third stage and the two
    flags (\csource{find_points.c}{3123}{3124}, \csource{find_points.c}{3129}{3129},
    \csource{rp-private.h}{577}{578}, \csource{rp-private.h}{583}{583}), gives
    \csource{sift.c}{405}{431} to Definition~\ref{def:stage3}, and also uses that definition; and
    Lemma~\ref{lem:sift_eq_checkDen} keeps \verb|Ratpoints.Sieve.toList_run_sift_eq_checkDen| with
    its new hypotheses, resting on Lemma~\ref{lem:staged}; the list of deviations of
    Chapter~\ref{chap:words} no longer names the stages, the binary $\gcd$ and the products as later
    work, nor a scan that reads a bit array in place, and the reduction of the row indices is left
    to stage~F3 alone;
  \item the run over the denominators (Definition~\ref{def:sieveRun}) takes the cut points, with
    \csource{find_points.c}{3861}{3861} and \csource{find_points.c}{3911}{3911}, makes the rows of
    the first $s_2$ entries and hands the rest to the third stage; its contract and the reduction of
    the run to the reference (Theorem~\ref{thm:run_reference}) need $b_{\mathrm{low}} \ge 1$ and the
    entries from $s_2$ on to be prime entries, and so do the passages on the run after
    Lemma~\ref{lem:sift_eq_checkDen} and in Section~\ref{sec:shapes}; the search
    (Definition~\ref{def:searchPoints}) passes the cut points of its parameters and reads the shape
    of the entries of the third stage off Lemma~\ref{lem:sievingInfo};
  \item the Jacobi test by tables (Definition~\ref{def:jacobiTables}) reduces with the reciprocal
    and uses Definition~\ref{def:mulMod}, its contract needs the primes of the table, its links to
    \verb|RP_MULDIV| move to Definition~\ref{def:mulMod}, and it holds the fact
    $\mathrm{mulModLimit} < 2^{32}$, which the files section of Chapter~\ref{chap:denominators}
    names among the facts that the proofs use; the deviation of Chapter~\ref{chap:denominators} on
    the divisions for \verb|RP_MULDIV| and \verb|RP_MULMOD| keeps only the reduction of the first
    word of a block modulo the forbidden divisors (stage~I);
  \item the sieve entry of a modulus (Lemma~\ref{lem:sieveEntry_admits}) and what
    \verb|sieving_info| computes (Lemma~\ref{lem:sievingInfo}) state the shape of the entries of the
    primes of the table and that the entries are those of the moduli in their order; the deviations
    that send the reciprocal of a sieve entry to stages~F2 and~F3 (in Chapter~\ref{chap:setup} the
    item on \verb|examine_prime|, in Chapter~\ref{chap:tables} the fields of the sieve entry of a
    prime and the last item) say that it is computed where it is read
    (Definition~\ref{def:mulMod}), the bias, \verb|rbainv| and \verb|dinv| serving the packing
    (stage~F3); the reciprocals left to stage~F2 in Chapter~\ref{chap:tables} are those of
    Definition~\ref{def:mulMod}, read by the third stage, the first two stages reading them in
    stage~F3, and the Jacobi test computes those of its primes in the same way where it reads them;
    the deviations of Chapter~\ref{chap:tables} on the primes beyond those examined, which ratpoints
    examines for its first two stages too when \verb|-p| is absent, and on the primes that \verb|-P|
    takes after the first $q$, which the parameters no longer take under \verb|-p|;
  \item the files sections of Chapters~\ref{chap:words} to~\ref{chap:denominators}, which name
    \verb|Ratpoints/Proof/Sift.lean| for the contracts from the walk on; the record of
    \verb|find_points_work_1| in Chapter~\ref{chap:setup}, which then lacks only what belongs to
    stages~F3, F4, H and~I; the passages that say that every modulus is applied in the first stage,
    that the cut points are not read, and that a run without options has a single stage.
\end{itemize}

\section{Files, names and order of work}\label{sec:stagesFiles}

The constants $\mathrm{stage3Limit}$, $\mathrm{useLongInPhase2}$, $\mathrm{defaultSp1}$ and
$\mathrm{sp2Extra}$, and the fact $\mathrm{mulModLimit} < 2^{32}$, go into
\verb|Ratpoints/Model/Constants.lean|, each citing its macro; the reciprocal and the reductions into
\verb|Ratpoints/Model/Tables.lean|; the test for common factors, the third stage, the second stage
and the staged sieve into \verb|Ratpoints/Model/Sieve.lean|; the run with the cut points and the
Jacobi test into \verb|Ratpoints/Model/Denominators.lean|; the search into
\verb|Ratpoints/Model/FindPoints.lean|; and the parameters into \verb|Ratpoints/Model/Params.lean|.
The contracts of the reductions, of the test for common factors, of the third stage with the test of
a candidate and of the second stage go into a new file without Mathlib,
\verb|Ratpoints/Proof/Stages.lean|; the corollary of the third stage for the tables of the reference
into \verb|Ratpoints/Proof/Tables.lean|; the contracts from the walk on (the walk, a bit array
through the second stage, the scan, one piece and one denominator) move from
\verb|Ratpoints/Proof/Words.lean| to a new file \verb|Ratpoints/Proof/Sift.lean|, which builds on
both, so that \verb|Words.lean| stays below the limit of $1500$ lines; the Jacobi test stays in
\verb|Ratpoints/Proof/JacobiTables.lean|, whose users in \verb|Ratpoints/Proof/Shapes.lean| follow;
the run in \verb|Ratpoints/Proof/Denominators.lean|, the sieve entries in
\verb|Ratpoints/Proof/Moduli.lean| and \verb|Ratpoints/Proof/FindPoints.lean|, and the parameters in
\verb|Ratpoints/Proof/Params.lean|. Every new module joins its part root and the listing
\verb|Ratpoints.lean|. One commit per item, each passing the project's checks, with the nodes of the
earlier chapters that it changes and those whose C lines it takes
(Section~\ref{sec:stagesEarlier}):
\begin{enumerate}
  \item[F2.1] The constant $\mathrm{stage3Limit}$ with its fact, the fact
    $\mathrm{mulModLimit} < 2^{32}$, the reciprocal and the reductions (Definition~\ref{def:mulMod},
    Lemma~\ref{lem:mulMod}), with \verb|#guard| examples against the remainder and the quotient over
    the table; nothing calls them yet, and the quotient shares the lines of \verb|RP_MULDIV| with
    the Jacobi test until F2.2.
  \item[F2.2] The Jacobi test by the quotient with the reciprocal, its contract with the primes of
    the table; the node of the Jacobi test, whose links to \verb|RP_MULDIV| pass to
    Definition~\ref{def:mulMod}.
  \item[F2.3] The test for common factors (Definition~\ref{def:relprime}, Lemma~\ref{lem:relprime})
    with its range lemma and \verb|#guard| examples against the $\gcd$; nothing calls it yet, and
    the check of a numerator gives it its lines (\csource{sift.c}{318}{343}); the node of
    Definition~\ref{def:checkNum}.
  \item[F2.4] The third stage and the test of a candidate (Definitions~\ref{def:stage3}
    and~\ref{def:accepted}, Lemmas~\ref{lem:stage3} and~\ref{lem:accepted}), with \verb|#guard|
    examples in both branches of the bias and above $\mathrm{mulModLimit}$; nothing calls them yet,
    and they take their lines from the first stage (\csource{sift.c}{291}{312}), from one
    denominator (\csource{sift.c}{405}{431}) and from the check of a numerator, which keeps none;
    the nodes of Definitions~\ref{def:andRows}, \ref{def:checkNum} and~\ref{def:sift}.
  \item[F2.5] The second stage (Definition~\ref{def:secondStage}, Lemma~\ref{lem:secondStage}), with
    the constant $\mathrm{useLongInPhase2}$; nothing calls them yet, and the first stage gives them
    \csource{sift.c}{1083}{1086} and \csource{sift.c}{1174}{1202}, the scan line~1297; the nodes of
    Definitions~\ref{def:andRows} and~\ref{def:scanPiece}.
  \item[F2.6] The staged sieve: a bit array through the second stage and the extraction
    (Definition~\ref{def:siftBitArray}, Lemma~\ref{lem:siftBitArray}), then the walk, the scan, the
    piece and one denominator with the test of a candidate, the second stage, the end of the first
    stage and the entries of the third, and the thunks (Lemma~\ref{lem:staged}); the new
    \verb|Ratpoints/Proof/Sift.lean|; the run passes all the entries to the first stage for now, and
    its contract gains only $b_{\mathrm{low}} \ge 1$; the nodes of Chapter~\ref{chap:words} and of
    the run.
  \item[F2.7] The run with the cut points, the search passing those of its parameters, the shape of
    the entries in the contracts of the sieve entry of a modulus and of \verb|sieving_info|; the
    theorems with their statements unchanged; the nodes of the run, of the search and of the sieve
    entries.
  \item[F2.8] The parameters of the command line by the counting rules of
    Section~\ref{sec:cuts}, with the constants $\mathrm{defaultSp1}$ and $\mathrm{sp2Extra}$ and
    \verb|#guard| examples of the stage sizes, confirmed by \verb|ratpoints -v| with \verb|-U 0|
    where \verb|-N| is absent; comparator, whose run covers the changes of F2.2 to F2.7 as well; the
    node of the parameters.
  \item[F2.9] The exclusions; the option sets of the differential tests; the build with other values
    of the constants, a Boolean among them and $\mathrm{mulModLimit}$ within its fact; the passages
    of the earlier chapters that F2 makes untrue; and the description of the repository and
    \verb|formalization.yaml|.
\end{enumerate}

\section{The evidence}\label{sec:stages-evidence}

Checks made while this chapter was written:
\begin{itemize}
  \item a reading of the C code of the stages at the pin, with a second reader who checked it: the
    stages and the caches are a data refinement of the unstaged sieve with its check, the numerators
    that reach the exact check and their order are the same, on 2300 runs of instrumented builds
    against a transcription; \verb|RP_MULMOD| and \verb|RP_MULDIV| exact for $u \delta < 2^{64}$,
    where $m p = 2^{64} + \delta$ with $\delta < p$, a bound that is sharp, and so for
    $u p < 2^{64}$; every \verb|long| of the stages in range for $H < 2^{63}$; the counting rules of
    the cut points, with five differences of the earlier parameters from them, which the parameters
    of Section~\ref{sec:cuts} remove;
  \item a sketch of the model of this chapter in Lean, whose \verb|#guard| examples give values of
    the C program; it agrees with \verb|ratpoints -q| on the 3219 runs of the differential tests,
    every point list in its order (in 22 seconds, against 35 before, since the second and the third
    stage now do part of the work), on the 13 runs at heights up to 100000 (in 116 seconds, against
    254), and on 1200 runs on 75 curves at heights up to 20000, without options and with fifteen
    option sets of \verb|-n|, \verb|-N|, \verb|-P| and \verb|-p|; where few moduli come before the
    test of a candidate, the model is much slower than ratpoints, up to 1209 seconds against 1.7 for
    \verb|-n 3 -N 5 -P 2| at the height 20000, since most numerators then reach the test and the
    reduction with the bias computes products above $2^{63}$ as big numbers; in a Lean test, the
    staged sieve gives the output of all the moduli in the first stage on 320 runs with random cut
    points, the test for common factors the $\gcd$ on 20000 pairs, the third stage after its table
    the rows of its primes on 4929 numerators (in both branches of the bias and above
    $\mathrm{mulModLimit}$), and the reductions the remainder and the quotient on 20000 arguments;
    the cut points of the parameters against those that \verb|ratpoints -v| reports on 1440 runs
    with \verb|-n| given and \verb|-U 0| where \verb|-N| is absent, the third stage compared where
    \verb|-P| is given: ratpoints has more moduli than the parameters in none, fewer in 275, each
    where its pool runs short; and the model built with other values of the constants of F2 (the
    arm of the second stage word by word, with all guards; and, with the guards left out, a third
    stage always without the bias with the cut points $53$ and $+0$, or a lower bound of the bias
    with $0$ and $+40$) gives the output of \verb|ratpoints -q| on the differential tests;
  \item two reviews of the sketch against the C code, each finding checked by a second reader: of
    the stages and the check, 26 findings, 24 confirmed or confirmed in part, several of them one
    finding of two readers, none on what the model computes (107000 cases of the arithmetic against
    the C code, none different); of the parameters, 16 findings, all on the text and the ranges; the
    rule bounds the counts of \verb|ratpoints -v| from above in every run where the options fix them
    (22080 runs: equal in 16912, below in 5168, where its pool runs short, above in none), and the
    well-formedness holds for 30 million choices of the options; the ranges and the deviations of
    this chapter are the corrected ones;
  \item a prototype of the proofs on the sketch, in a copy of the repository, in three parts in
    sequence: every contract and lemma of this chapter proved, with the standard axioms only, every
    proof of at most $20$ lines, the statements that comparator certifies and the comparator
    solution unchanged in their text, the build with other values of the constants, both arms of the
    second stage among them, and the differential tests of its program; it changed nothing in the
    model but the fact $\mathrm{mulModLimit} < 2^{32}$, which the exactness of the reductions of the
    denominators needs, and the entries of the third stage need of the tables only the inverses of
    the nonzero residues and the bound $p \le 2^{32}$;
  \item a review of this chapter by four readers, each finding checked by a second reader: 64
    findings, 44 confirmed, 19 in part and 1 refuted; the prototype then gave the hypotheses of the
    third stage as prime entries; the rest was applied to this chapter: the order of work, in which
    each item takes the C lines of its components from the earlier nodes in its own commit and a bit
    array through the second stage comes with the walk that it calls, the positive denominators of
    the run, the changes of the earlier chapters, the C links, the uses of the nodes, and the claims
    of the evidence.
\end{itemize}
